\chapter{Documentación de Extracción}
Este anexo presenta una descripción general de los \textit{scripts} que conforman el componente de extracción. 
Las instrucciones de uso correspondientes, así como los detalles específicos del formato y la estructura de los 
archivos generados durante la ejecución, se detallan en el Anexo~\ref{anx:manual_extraccion}.

\section{Requisitos}
El componente requiere un entorno de ejecución basado en Python 3.12.10 y las siguientes dependencias de software:

\begin{verbatim}
    filelock==3.32.0
    fsspec==2026.6.0
    Jinja2==3.1.6
    MarkupSafe==3.0.3
    mpmath==1.3.0
    networkx==3.6.1
    numpy==2.5.1
    pillow==12.3.0
    setuptools==83.0.0
    sympy==1.13.1
    torch==2.6.0+cu124
    torchvision==0.21.0+cu124
    typing_extensions==4.16.0
\end{verbatim}
Adicionalmente, se requiere el módulo de enlace para Python de la biblioteca \textit{OpenVDB} (\texttt{pyopenvdb}) 
compilado para la versión específica del intérprete en uso. El procedimiento detallado de compilación e instalación 
se documenta en el Anexo~\ref{anx:compilacion_openvdb}.
\medskip

Por último, el componente integra la implementación en PyTorch de Tiny YOLOv2 desarrollada por \textcite{yolov2_implementation}.

\section{Organización del contenido}
El componente de extracción se compone de cuatro \textit{scripts} principales de Python:
\begin{itemize}
    \item \texttt{ejecucion.py}: Módulo encargado de instanciar la arquitectura Tiny YOLOv2, transferir los pesos preentrenados, acondicionar el tensor de entrada y 
    gestionar la inferencia del modelo.
    \item \texttt{extractor.py}: Administra la suscripción de funciones de intercepción (\textit{forward hooks}) para capturar las activaciones intermedias durante 
    la propagación hacia adelante, recopilando además la información topológica y metadatos de las capas.
    \item \texttt{exportar.py}: Recupera los tensores residentes en memoria y metadatos y gestiona su persistencia y serialización en los formatos estructurados (\textit{OpenVDB} 
    y secuencias de imágenes).
    \item \texttt{main.py}: Actúa como punto de entrada y orquestador central del flujo de trabajo (\textit{pipeline}), exponiendo la interfaz de línea de comandos.
\end{itemize}

Asimismo, el módulo se apoya en los siguientes componentes auxiliares:
\begin{itemize}
    \item \texttt{TinyYoloV2/}: Directorio que contiene la implementación en PyTorch de la arquitectura Tiny YOLOv2 basada en \parencite{yolov2_implementation}.
    \item \texttt{utils.py}: Provee funciones de soporte para la verificación, saneamiento y creación recursiva de rutas y directorios en el sistema de archivos.
    \item \texttt{benchmark.py}: Módulo complementario encargado de coordinar la ejecución por lotes y recopilar métricas comparativas de rendimiento temporal 
    y precisión.
\end{itemize}

La disposición jerárquica del paquete se estructura de la siguiente manera:
\begin{figure}[!htbp]
\dirtree{%
.1 componente-extraccion/.
.2 benchmark.py.
.2 ejecucion.py.
.2 exportar.py.
.2 extractor.py.
.2 main.py.
.2 utils.py.
.2 TinyYoloV2/.
.3 model.py.
.3 utils.py.
}
\caption{Estructura del directorio del componente de extracción.}
\label{fig:arbol_componente_extraccion}
\end{figure}

\section{Scripts}

\subsection{\texttt{ejecucion.py}}
El \textit{script} \texttt{ejecucion.py} administra la inicialización, parametrización y ejecución de la inferencia sobre la arquitectura neuronal a 
partir de una imagen de entrada. Toda la lógica operativa se encapsula en la clase \texttt{ModelExecutor} para garantizar modularidad e independencia funcional.
\medskip

Esta clase preprocesa la imagen de entrada convirtiéndola en un tensor normalizado con las dimensiones esperadas por la red ($1 \times 3 \times 416 \times 416$). 
Posteriormente, instancia la arquitectura Tiny YOLOv2, transfiere los pesos preentrenados a la memoria y ubica el modelo en el dispositivo de cómputo seleccionado 
por PyTorch. Al completar la inferencia, produce una lista con los vectores de las cajas delimitadoras resultantes 
(filtradas mediante umbrales de confianza y supresión de no máximos, NMS) junto con una previsualización gráfica que sobrepone las detecciones a la imagen original.

\subsubsection*{Atributos de \texttt{ModelExecutor}}
\begin{itemize}
    \item \texttt{device}: Dispositivo asignado para la ejecución en PyTorch (\texttt{torch.device}). Selecciona de forma prioritaria un acelerador 
    compatible con CUDA o conmuta a CPU en caso de no disponer de uno.
    \item \texttt{weight\_path}: Ruta del archivo binario con los pesos preentrenados del modelo (formato \texttt{.weights}).
    \item \texttt{conf\_threshold}: Umbral de confianza (\textit{objectness threshold}) empleado para discriminar predicciones candidatas.
    \item \texttt{nms\_threshold}: Umbral de intersección sobre unión (IoU) utilizado por el algoritmo de supresión de no máximos.
    \item \texttt{input\_path}: Ruta del archivo de la imagen suministrada como entrada.
    \item \texttt{input\_image}: Imagen de entrada cargada en memoria (matriz RGB / objeto de imagen).
    \item \texttt{input\_tensor}: Tensor de entrada preprocesado y transferido al dispositivo de cómputo.
    \item \texttt{boxes}: Colección de predicciones y cajas delimitadoras calculadas por el modelo.
\end{itemize}

\subsubsection*{Métodos de \texttt{ModelExecutor}}
\begin{itemize}
    \item \texttt{\_\_init\_\_(input\_path)}: Constructor de la clase; inicializa los parámetros de inferencia y asigna la ruta de la imagen a procesar.
    \item \texttt{configure\_input\_image()}: Ejecuta la secuencia de acondicionamiento de entrada invocando la lectura de la imagen y la construcción del tensor de activación inicial.
    \item \texttt{\_load\_image()}: Lee la imagen desde el sistema de archivos a partir de \texttt{input\_path}.
    \item \texttt{\_load\_input\_tensor()}: Convierte y normaliza la imagen leída en un tensor compatible con el grafo de ejecución.
    \item \texttt{\_load\_model()}: Instancia el modelo Tiny YOLOv2, transfiere los coeficientes de los pesos y ubica los tensores en el dispositivo configurado.
    \item \texttt{execute()}: Ejecuta la inferencia del modelo y aplica la decodificación vectorial junto con el filtrado por NMS.
    \item \texttt{clear()}: Libera los recursos de memoria asignados a las imágenes de entrada/salida y a las estructuras de cajas generadas.
    \item \texttt{get\_boxes()}: Método de acceso (\textit{getter}) que retorna la lista estructurada de cajas delimitadoras calculadas.
\end{itemize}

\subsection{\texttt{extractor.py}}
Este módulo implementa los mecanismos de introspección para capturar las activaciones intermedias y la topología de la red Tiny YOLOv2. 
Dicha funcionalidad se implementa en la clase \texttt{ModelExtractor}.
\medskip

Tomando como dependencia una instancia de \texttt{ModelExecutor}, la clase registra funciones de intercepción (\textit{forward hooks}) 
en las capas seleccionadas para capturar los tensores residentes en memoria durante el paso hacia adelante, además de recopilar metadatos 
analíticos de la arquitectura. Adicionalmente, implementa directivas de filtrado para delimitar la extracción mediante listas de inclusión y 
exclusión explícitas. Los datos extraídos se consolidan en una estructura jerárquica indexada por capa.

\subsubsection*{Atributos de \texttt{ModelExtractor}}
\begin{itemize}
    \item \texttt{input\_image}: Referencia a la imagen de entrada leída por el ejecutor.
    \item \texttt{output\_image}: Imagen de salida con las anotaciones de las cajas superpuestas.
    \item \texttt{boxes}: Lista de vectores numéricos correspondientes a las predicciones filtradas.
    \item \texttt{layers}: Diccionario con los tensores de activación extraídos, indexados por el identificador único de cada capa.
    \item \texttt{layers\_data}: Diccionario de metadatos estructurales (tipo de capa, dimensión del núcleo convolucional, paso o \textit{stride}, 
    relleno o \textit{padding}, etc.), indexado por el nombre de la capa.
    \item \texttt{model\_exe}: Instancia referenciada de la clase \texttt{ModelExecutor}.
    \item \texttt{model}: Referencia al módulo de red neuronal instanciado en el ejecutor.
    \item \texttt{classes}: Lista con las etiquetas categóricas del conjunto de datos.
    \item \texttt{capture\_layers}: Lista de clases u operadores (e.g., \textit{Conv2d}) cuyos tensores deben ser interceptados sistemáticamente.
    \item \texttt{include}: Lista de nombres canónicos de capas específicas a incluir de forma forzada, con precedencia sobre \texttt{capture\_layers}.
    \item \texttt{exclude}: Lista de nombres canónicos de capas a omitir durante la extracción, con precedencia sobre \texttt{capture\_layers} e \texttt{include}.
\end{itemize}

\subsubsection*{Métodos de \texttt{ModelExtractor}}
\begin{itemize}
    \item \texttt{\_\_init\_\_(model\_exe, \dots)}: Constructor de la clase; recibe la instancia de \texttt{ModelExecutor} e inicializa los registros de captura.
    \item \texttt{hook\_fn(module, input, output)}: Función de intercepción ejecutada durante el paso hacia adelante (\textit{forward step}) para registrar y duplicar los tensores de activación.
    \item \texttt{\_get\_layer\_data()}: Inspecciona las propiedades constructivas de los módulos y suscribe formalmente los \textit{forward hooks}.
    \item \texttt{\_get\_layer\_connections()}: Reconstruye el grafo computacional para derivar las dependencias emisor-receptor (padre-hijo) entre capas consecutivas.
    \item \texttt{\_node\_name(node)}: Función auxiliar de \texttt{\_get\_layer\_connections()} encargada de normalizar la denominación de los nodos dentro del grafo.
    \item \texttt{print\_info()}: Despliega en consola el esquema topológico de la red, listando el identificador de cada módulo y su tipo de operación asociada.
    \item \texttt{extract()}: Dispara la inferencia en el ejecutor y orquesta la recolección de tensores y metadatos derivados.
    \item \texttt{clear()}: Vacía las colecciones y diccionarios de tensores intermedios liberando la memoria asignada.
\end{itemize}

\subsection{\texttt{exportar.py}}
Este módulo gestiona la persistencia de las activaciones y metadatos capturados en el sistema de almacenamiento. Soporta la 
serialización a cuadrículas volumétricas dispersas en formato \textit{OpenVDB} (\texttt{.vdb}), así como la descomposición 
en imágenes bidimensionales rasterizadas (\texttt{.jpg}, \texttt{.png} y \texttt{.tiff}) para análisis comparativo y depuración. 
Dicho comportamiento se abstrae en la clase \texttt{ModelExporter}.
\medskip

Durante el proceso, el módulo genera un directorio de destino donde preserva copias de las imágenes de entrada y salida, 
serializa la estructura \texttt{layers\_data} en el archivo \texttt{manifest.json} y exporta las predicciones a la cuadrícula 
volumétrica \texttt{boxes.vdb}. En este último archivo, cada predicción es mapeada como un vector discreto a lo largo de los 
ejes espaciales, consolidando todas las detecciones en una única estructura que posteriormente es reconstruida por las herramientas de visualización.

\subsubsection*{Atributos de \texttt{ModelExporter}}
\begin{itemize}
    \item \texttt{supported\_formats}: Tupla o conjunto que define las extensiones y formatos de persistencia admitidos por el exportador.
    \item \texttt{dry\_run}: Indicador booleano de depuración; ejecuta la totalidad del flujo computacional y de conversión omitiendo la escritura física en disco.
    \item \texttt{formats}: Lista de formatos de exportación seleccionados para la corrida en curso.
    \item \texttt{info}: Indicador booleano para habilitar la salida detallada (\textit{verbose}) de telemetría durante la serialización.
    \item \texttt{model\_extractor}: Instancia referenciada de la clase \texttt{ModelExtractor}.
\end{itemize}

\subsubsection*{Métodos de \texttt{ModelExporter}}
\begin{itemize}
    \item \texttt{\_\_init\_\_(model\_extractor)}: Constructor; vincula la fuente de tensores y metadatos 
    a partir de una instancia de \texttt{ModelExtractor}.
    \item \texttt{export(output\_path)}: Método principal de exportación de los datos, guardando los 
    archivos resultantes en \texttt{output\_path}.
    \item \texttt{\_save\_io\_images(output\_dir)}: Escribe en el directorio de destino los archivos 
    gráficos de entrada y salida anotada.
    \item \texttt{\_convert\_boxes(output\_dir)}: Codifica matricialmente las cajas delimitadoras resultantes 
    y las serializa en el archivo \texttt{boxes.vdb}.
    \item \texttt{\_convert\_tensor(layer\_name)}: Transforma el tensor en memoria de la capa especificada en 
    una cuadrícula volumétrica de \textit{OpenVDB}.
    \item \texttt{\_tensor\_to\_images(layer\_name, ext)}: Procesa el tensor de activación descomponiéndolo a 
    lo largo de sus canales en una secuencia de matrices bidimensionales.
    \item \texttt{\_save\_tensor\_as\_grid(layer\_name, output\_dir)}: Serializa la cuadrícula \textit{OpenVDB} 
    de una capa determinada escribiendo el archivo \texttt{.vdb} en el almacenamiento secundario.
    \item \texttt{\_save\_tensor\_as\_image\_seq(layer\_name, output\_dir, ext)}: Exporta los canales del tensor 
    omo una secuencia indexada de imágenes rasterizadas en la ruta indicada.
    \item \texttt{\_get\_exported\_volume\_as\_numpy32(layer\_name, ext)}: Lee un volumen previamente serializado 
    y lo reconstruye en un arreglo \textit{NumPy} de precisión simple (\texttt{float32}) para validación de error.
    \item \texttt{\_get\_vdb\_grid\_as\_numpy32(grid)}: Convierte los vóxeles de una cuadrícula de \textit{OpenVDB} 
    en un arreglo matricial de tipo \texttt{float32}.
    \item \texttt{\_get\_images\_as\_numpy32(images)}: Reconstruye una secuencia de imágenes leídas de disco en una 
    estructura tensorial equivalente de \textit{NumPy} en precisión simple (\texttt{float32}).
\end{itemize}

\subsection{\texttt{main.py}}
El \textit{script} \texttt{main.py} actúa como orquestador y punto de entrada central del flujo de trabajo 
(\textit{pipeline}). Integra y coordina el ciclo de vida de las clases \texttt{ModelExecutor}, \texttt{ModelExtractor} y 
\texttt{ModelExporter}. Asimismo, define y analiza los argumentos de la interfaz de línea de comandos (CLI), permitiendo 
parametrizar dinámicamente las rutas de entrada/salida, los formatos de destino, los umbrales de detección y los filtros 
de capas a procesar.

\subsection{\texttt{utils.py}}
El módulo \texttt{utils.py} agrupa funciones auxiliares para la verificación y saneamiento de rutas en el sistema de archivos:
\begin{itemize}
    \item \texttt{directory\_validation(path)}: Comprueba la validez de una ruta de directorio y asegura su existencia en disco creándola 
    de forma recursiva si no se encuentra presente.
    \item \texttt{file\_validation(path)}: Valida formalmente la existencia física y la integridad de acceso del archivo especificado.
\end{itemize}

\subsection{Subdirectorio \texttt{TinyYoloV2/}}
Este directorio contiene la adaptación en PyTorch de la arquitectura Tiny YOLOv2 basada en la implementación de \parencite{yolov2_implementation}. 
Los dos componentes primarios de este paquete son:
\begin{itemize}
    \item \texttt{model.py}: Define la clase del modelo neuronal, su topología convolucional y el método de propagación hacia adelante (\textit{forward pass}).
    \item \texttt{utils.py}: Provee rutinas complementarias para la decodificación matricial de las predicciones, el filtrado mediante supresión 
    de no máximos (NMS) y el renderizado gráfico de las cajas delimitadoras sobre la imagen de salida.
\end{itemize}

\subsection{\texttt{benchmark.py}}
Este módulo orquesta la evaluación de rendimiento computacional ejecutando el flujo de trabajo por lotes (\textit{batch processing}) 
sobre un conjunto de imágenes. Expone las siguientes variables globales de configuración:

\begin{itemize}
    \item \texttt{INPUT\_DIR}: Ruta del directorio que contiene el conjunto de imágenes de prueba a procesar.
    \item \texttt{OUTPUT\_DIR}: Directorio raíz de destino para los artefactos extraídos. En su interior, los resultados se 
    estructuran en subdirectorios nombrados según el identificador de la imagen de entrada correspondiente.
    \item \texttt{CAPTURE\_LAYERS}: Lista con los tipos de capas u operadores del modelo a interceptar.
    \item \texttt{EXECUTION\_COUNT}: Cantidad total de imágenes a procesar en la corrida de prueba.
    \item \texttt{FORMAT}: Formato de serialización y exportación de las cuadrículas tensoriales.
\end{itemize}

Asimismo, inicializa a nivel global las instancias de los componentes del flujo:
\begin{itemize}
    \item \texttt{executor}: Instancia de \texttt{ModelExecutor} encargada de la inferencia del modelo.
    \item \texttt{extractor}: Instancia de \texttt{ModelExtractor} configurada para la captura de tensores intermedios.
    \item \texttt{exporter}: Instancia de \texttt{ModelExporter} responsable de la persistencia de los datos.
\end{itemize}
El \textit{script} mantiene una única instancia activa del ejecutor del modelo compartida a lo largo de toda la secuencia, optimizando el tiempo 
global de procesamiento al evitar la recarga redundante de los pesos preentrenados y la reinicialización de grafos en memoria entre iteraciones consecutivas.
\medskip

El módulo implementa las siguientes funciones de soporte y control:
\begin{itemize}
    \item \texttt{progress\_bar(count, total)}: Rutina auxiliar que despliega una barra de progreso interactiva en la terminal para monitorizar 
    el avance del procesamiento por lotes.
    \item \texttt{run(input\_file, output\_path)}: Ejecuta el ciclo integral de inferencia, extracción y serialización para una muestra individual. 
    Registra la telemetría temporal de la ejecución y los metadatos de los archivos generados.
    \item \texttt{output\_tree\_info(path)}: Analiza de forma recursiva la jerarquía del directorio generado, calculando el tamaño en disco 
    (en bytes) y la cantidad de archivos contenidos en cada subcarpeta.
    \item \texttt{calculate\_precision\_loss(ext)}: Cuantifica la pérdida de precisión numérica generada por el proceso de exportación para 
    una extensión dada. La discrepancia se evalúa contrastando directamente el tensor original en memoria contra la estructura intermedia 
    procesada por el exportador antes de su escritura física, evitando sobrecostos por operaciones de lectura en disco.
    \item \texttt{batch\_run(batch\_count, batch\_output\_path)}: Orquesta la ejecución secuencial del procesamiento por lotes sobre las 
    imágenes del directorio de entrada. Consolida los artefactos y el reporte agregado de telemetría bajo un directorio único de salida y 
    despliega en consola las métricas estadísticas promedio.
\end{itemize}

Además de los artefactos de extracción correspondientes a cada muestra, el módulo genera un archivo de balance \texttt{manifest.json} en 
la raíz del directorio de salida con la telemetría individual y las métricas agregadas del lote, indexadas por el nombre del archivo 
de entrada correspondiente.

Las métricas registradas de forma individual por cada ejecución son:
\begin{itemize}
    \item \texttt{Input File}: Nombre o ruta del archivo de imagen suministrado a la red.
    \item \texttt{Output Directory}: Directorio de destino asignado a la muestra procesada.
    \item \texttt{Runtime}: Tiempo total de ejecución del flujo individual (en segundos), abarcando la inferencia del modelo, 
    la captura de tensores y la exportación a disco.
    \item \texttt{Total Filesize}: Tamaño acumulado de los artefactos generados (en bytes), excluyendo las cuadrículas de cajas 
    delimitadoras y las imágenes de entrada y salida.
    \item \texttt{Total Filecount}: Cantidad neta de archivos producidos, aplicando los mismos criterios de exclusión de metadatos 
    y previsualizaciones.
    \item \texttt{Run Parameters}: Parámetros de configuración suministrados al flujo, tales como \texttt{--include}, 
    \texttt{--exclude}, \texttt{--conf} y \texttt{--nms}, entre otros.
    \item \texttt{Precision Error}: Pérdida de precisión numérica estimada entre los tensores residentes en memoria y los 
    datos exportados (calculada mediante el error absoluto medio o MAE). Las métricas se agrupan por formato e incluyen 
    los valores promedio, mínimo y máximo, además del error específico cuantificado por cada capa intervenida.
    \item \texttt{Files}: Estructura jerárquica de los archivos y carpetas generados, documentando el tamaño en disco de 
    cada elemento (en bytes) y la cantidad de archivos contenidos por subdirectorio.
\end{itemize}

Por su parte, las métricas globales agregadas para el lote completo comprenden:
\begin{itemize}
    \item \texttt{Runs}: Cantidad total de ejecuciones individuales procesadas.
    \item \texttt{Min Runtime}: Mínimo tiempo de ejecución individual registrado en el lote (en segundos).
    \item \texttt{Max Runtime}: Máximo tiempo de ejecución individual registrado en el lote (en segundos).
    \item \texttt{Avg Runtime}: Media aritmética del tiempo de ejecución por muestra (en segundos).
    \item \texttt{Total Runtime}: Tiempo acumulado de procesamiento (en segundos), computando exclusivamente las 
    fases de inferencia, captura de tensores y exportación en disco.
    \item \texttt{Total Filesize}: Espacio total ocupado en disco por los artefactos generados (en bytes).
    \item \texttt{Total Filecount}: Cantidad total de archivos generados a lo largo del lote.
    \item \texttt{MAE}: Pérdida de precisión numérica media estimada por tipo de formato o extensión 
    (calculada mediante el error absoluto medio).
\end{itemize}